Paper~7~\cite{paper7online} found that lag-1 online recalibration of
HygienePrio harms performance at $K = 200$ (cell-mean recovery ratio
$\rho = -0.66$). Paper~8~\cite{paper8multi} falsified the natural
fix --- multi-window smoothing makes the hazard worse, not better.
Paper~8 explicitly named adaptive change-point detection as the
remaining candidate. This paper evaluates the simplest such
detector: at each window, monitor the calibration-target P@50 shift
between consecutive past windows and switch to fixed weights when
the magnitude crosses a pre-registered threshold of $0.05$.

Five strategies (\textsf{fixed}, \textsf{lag1}, \textsf{adaptive},
\textsf{gated}, \textsf{offline}) were evaluated on a pre-registered
2{,}250-row sweep ($K \in \{50, 100, 200\}$, $\lambda = 3$, six
windows, 25 evaluation seeds). The pre-registration locked four
hypotheses; the outcome is mixed.

\textbf{H1 supported:} At $K = 200$, \textsf{adaptive} is the first
procedure in the Paper 7--8--10 sequence that does not harm relative
to \textsf{fixed} at any of the six windows. Mean delta over the
five post-W1 windows is $+0.7$~pp, with the W2--W4 windows up to
$+3.4$~pp. The high-capacity hazard is finally avoided.

\textbf{H4 supported:} At $K = 200$, \textsf{adaptive} marginally
beats a static capacity-conditional gate (\textsf{gated}) by $+0.7$~pp
($0.273$ vs $0.266$, cell mean). The per-window detector adds value
over the static rule, but the margin is small.

\textbf{H2 rejected:} At $K = 50$, the detector fires in three of
the five post-W1 windows ($60\%$), and the adaptive cell mean is
$5.3$~pp \textit{below} lag-1's because every fire reverts to
\textsf{fixed} when lag-1 would have done better. Adaptive trades
moderate-capacity recovery for high-capacity safety.

\textbf{H3 rejected:} The firing fraction is not monotone in $K$
($60\%, 40\%, 80\%$ at $K \in \{50, 100, 200\}$; Spearman
$\rho = 0.5$). The detector picks up on per-window distributional
shifts that are not cleanly aligned with capacity.

The deployable conclusion is operational rather than algorithmic:
\textbf{adaptive switching is the first procedure that does not
break at high capacity, but it pays a moderate-capacity penalty
and is only marginally superior to a static capacity-conditional
gate}. A practitioner can prefer the simpler gate ($+0.7$~pp gap)
unless capacity is variable across deployment cycles. All results
bounded to the synthetic EEHDA evaluation context; external validation
on real fleet telemetry is required before any deployment recommendation.
